\section{Conclusion}
\label{sec:conclusion}

This paper set out to test whether \citet{andersen2025}'s OLS finding that the Scottish Child Payment reduced child material deprivation survives re-estimation with a doubly robust, machine-learning estimator, and whether that aggregate effect is driven by genuine change in specific, identifiable hardships. Using nine years of FRS/HBAI data and three independent estimation strategies, the answer to both is a qualified yes and no: the aggregate effect survives, though roughly halved once functional-form assumptions are relaxed ($-0.076$ OLS to $-0.028$ to $-0.045$ DML), and remains directionally consistent across every specification tested; but no single deprivation item carries a robust, correctable signal of its own, so the aggregate finding should be read as evidence of a real but diffuse improvement in material circumstances, not a traceable effect on any one hardship.

For policy design, this matters beyond Scotland. Wales's Cynnal pilot (Section~\ref{sec:introduction}) is, like SCP was in 2021, a modest weekly transfer aimed at the same population; this paper's evidence suggests such a transfer can plausibly move material deprivation, but any evaluation of Cynnal should not expect, and should not require, a detectable effect on any single measured hardship to conclude the transfer worked. The gap between this paper's doubly robust estimate and the Scottish Government's own income-poverty modelling (Section~\ref{sec:discussion}) is also a reminder that a payment's effect on income poverty and on material deprivation are related but distinct questions, both worth asking of any future UK child payment.

The clearest limitations -- an ITT estimand that dilutes the effect on actual recipients, an unresolved residual spillover risk, and an SCP effect not fully separable from Scotland's other child transfers (Section~\ref{sec:discussion}) -- are all directions for future work once further years of post-rollout data, and ideally administrative receipt data, become available.
